% ═══════════════════════════════════════════════════════════════════
% §7  Ground Source Heat Pump Boiler (Dynamic Model)
% Source: GroundSourceHeatPumpBoiler.py
% ═══════════════════════════════════════════════════════════════════
\section{Ground source heat pump boiler}
\label{sec:gshpb}

This section describes the dynamic GSHPB model in \texttt{GroundSourceHeatPumpBoiler.py}. The model shares the vapour-compression framework of the ASHPB (Section~\ref{sec:ashpb}) but replaces the air-source evaporator with a ground-coupled borehole heat exchanger.

\subsection{Assumptions}
\begin{enumerate}
  \item Ground temperature is uniform and constant at depth (undisturbed ground temperature $T_g$).
  \item The borehole heat exchanger is modelled via the finite line source (FLS) g-function.
  \item Fluid properties in the ground loop are evaluated at the mean brine temperature.
  \item The condenser model is identical to the ASHPB (submerged coil, $UA\cdot\Delta T$).
\end{enumerate}

\subsection{Ground-coupled evaporator}

The heat extraction from the ground:
\begin{equation}\label{eq:gshp-Qground}
  Q_\text{ground} = \frac{T_g - T_\text{brine,in}}{R_b + R_g(t)}
\end{equation}
where $R_b$ is the borehole thermal resistance and $R_g(t)$ is the time-dependent ground thermal resistance computed from the g-function (Eq.~\ref{eq:fls-g}):
\begin{equation}\label{eq:gshp-Rg}
  R_g(t) = \frac{g(t)}{2\pi k_s H}
\end{equation}
with $k_s$ the ground thermal conductivity and $H$ the borehole depth.

\subsection{Modified COP correlation}

The COP is based on a Carnot-modified formula using the ground-loop brine temperature:
\begin{equation}\label{eq:gshp-COP}
  \text{COP}_\text{GSHP} = \eta_\text{Carnot}\,\frac{T_\text{cond}}{T_\text{cond} - T_\text{brine}}
\end{equation}
where $\eta_\text{Carnot}$ accounts for real-cycle losses (typically 0.4--0.6)~\cite{sarbu2019}.

\subsection{Brine loop energy balance}

\begin{equation}\label{eq:gshp-brine}
  T_\text{brine,out} = T_\text{brine,in} - \frac{Q_\text{evap}}{\dot{m}_\text{brine}\,c_\text{brine}}
\end{equation}

\subsection{Exergy analysis}

The exergy balance structure is analogous to the ASHPB (Section~\ref{sec:ashpb}), with the outdoor unit replaced by the ground-coupled loop:
\begin{equation}\label{eq:gshp-Xground}
  X_\text{ground} = Q_\text{ground}\left(1 - \frac{T_0}{T_g}\right)
\end{equation}

The overall exergy efficiency:
\begin{equation}\label{eq:gshp-etaX}
  \eta_{X,\text{sys}} = \frac{X_\text{w,out} - X_\text{w,in}}{X_\text{cmp} + X_\text{pump,brine}}
\end{equation}
